\textbf{Phần 1. Trắc nghiệm lựa chọn}
\setcounterpageref{ex}{0}

\begin{ex} %[2D1N02]
Trên tia số, gốc của tia số là số nào?
\choice
{\True $0$}
{$1$}
{$10$}
{$100$}
\loigiai{
Gốc của tia số là số 0.\\
\textbf{Đáp án của câu hỏi là A}
}
\end{ex}

\begin{ex} %[2D1N02]
Số liền sau của số $39$ là
\choice
{$38$}
{\True $40$}
{$41$}
{$30$}
\loigiai{
$39 + 1 = 40$.\\
\textbf{Đáp án của câu hỏi là B}
}
\end{ex}

\begin{ex} %[2D1H02]
Số liền trước của số $70$ là
\choice
{$71$}
{$60$}
{\True $69$}
{$80$}
\loigiai{
$70 - 1 = 69$.\\
\textbf{Đáp án của câu hỏi là C}
}
\end{ex}

\begin{ex} %[2D1H02]
Số liền trước của số $1$ là
\choice
{\True $0$}
{$2$}
{$10$}
{Không có}
\loigiai{
$1 - 1 = 0$.\\
\textbf{Đáp án của câu hỏi là A}
}
\end{ex}

\begin{ex} %[2D1H02]
Trên tia số, số đứng liền bên phải số 45 là
\choice
{$44$}
{$47$}
{$50$}
{\True $46$}
\loigiai{
Đứng liền bên phải là số liền sau: $45 + 1 = 46$.\\
\textbf{Đáp án của câu hỏi là D}
}
\end{ex}

\begin{ex} %[2D1H02]
Trên tia số, số đứng liền bên trái số 80 là
\choice
{$81$}
{$70$}
{\True $79$}
{$90$}
\loigiai{
Đứng liền bên trái là số liền trước: $80 - 1 = 79$.\\
\textbf{Đáp án của câu hỏi là C}
}
\end{ex}

\begin{ex} %[2D1V02]
Số liền sau của số lớn nhất có hai chữ số là
\choice
{\True $100$}
{$98$}
{$99$}
{$101$}
\loigiai{
Số lớn nhất có 2 chữ số là 99. Số liền sau là $99 + 1 = 100$.\\
\textbf{Đáp án của câu hỏi là A}
}
\end{ex}

\begin{ex} %[2D1V02]
Số liền trước của số bé nhất có hai chữ số là
\choice
{$10$}
{$11$}
{\True $9$}
{$8$}
\loigiai{
Số bé nhất có 2 chữ số là 10. Số liền trước là $10 - 1 = 9$.\\
\textbf{Đáp án của câu hỏi là C}
}
\end{ex}

\begin{ex} %[2D1V02]
Ba số liên tiếp tăng dần là
\choice
{$18, 20, 19$}
{\True $28, 29, 30$}
{$40, 39, 38$}
{$50, 52, 54$}
\loigiai{
Ba số liên tiếp: $28, 29, 30$.\\
\textbf{Đáp án của câu hỏi là B}
}
\end{ex}

\begin{ex} %[2D1V02]
Tìm số ở giữa số $68$ và số $70$:
\choice
{$67$}
{$71$}
{$72$}
{\True $69$}
\loigiai{
Số ở giữa 68 và 70 là 69.\\
\textbf{Đáp án của câu hỏi là D}
}
\end{ex}

\begin{ex} %[2D1V02]
Có bao nhiêu số tự nhiên nằm giữa hai số $25$ và $29$?
\choice
{$2$ số}
{\True $3$ số ($26, 27, 28$)}
{$4$ số}
{$5$ số}
\loigiai{
Các số là $26, 27, 28$ (có 3 số).\\
\textbf{Đáp án của câu hỏi là B}
}
\end{ex}

\begin{ex} %[2D1C02]
Tìm một số biết rằng số liền trước của nó là số lớn nhất có 1 chữ số.
\choice
{$9$}
{$11$}
{$8$}
{\True $10$}
\loigiai{
Số lớn nhất có 1 chữ số là 9. Số cần tìm là $9 + 1 = 10$.\\
\textbf{Đáp án của câu hỏi là D}
}
\end{ex}

\textbf{Phần 2. Trắc nghiệm đúng sai}
\setcounterpageref{ex}{0}

\begin{ex} %[2D1H02]
Xét tính đúng sai về số liền trước và số liền sau:
\choiceTF
{\True Số liền sau hơn số liền trước 1 đơn vị.}
{\True Số liền trước của 50 là 49.}
{Số liền sau của 89 là 88.}
{\True Số liền sau của 89 là 90.}
\loigiai{
\begin{itemize}
    \item Liền sau hơn liền trước 1 đơn vị: Đúng.
    \item $50 - 1 = 49$: Đúng.
    \item $89 + 1 = 90 \neq 88$. Ý 3 Sai, Ý 4 Đúng.
\end{itemize}
\begin{itemize}
    \item a) \textbf{Đúng}.
    \item b) \textbf{Đúng}.
    \item c) \textbf{Sai}.
    \item d) \textbf{Đúng}.
\end{itemize}
}
\end{ex}

\begin{ex} %[2D1V02]
Cho số 35. Xét tính đúng sai:
\choiceTF
{\True Số liền trước của 35 là 34.}
{\True Số liền sau của 35 là 36.}
{Tổng của số liền trước và số liền sau của 35 là 60.}
{\True Tổng của số liền trước và số liền sau của 35 là 70 ($34 + 36 = 70$).}
\loigiai{
\begin{itemize}
    \item Liền trước là 34: Đúng.
    \item Liền sau là 36: Đúng.
    \item $34 + 36 = 70 \neq 60$. Ý 3 Sai, Ý 4 Đúng.
\end{itemize}
\begin{itemize}
    \item a) \textbf{Đúng}.
    \item b) \textbf{Đúng}.
    \item c) \textbf{Sai}.
    \item d) \textbf{Đúng}.
\end{itemize}
}
\end{ex}

\textbf{Phần 3. Trả lời ngắn}
\setcounterpageref{ex}{0}

\begin{ex} %[2D1H02]
Tìm số liền trước của số $85$.
\shortans{84}
\loigiai{
\textbf{Đáp số:} $84$\\
$85 - 1 = 84$.
}
\end{ex}

\begin{ex} %[2D1H02]
Tìm số liền sau của số $79$.
\shortans{80}
\loigiai{
\textbf{Đáp số:} $80$\\
$79 + 1 = 80$.
}
\end{ex}

\begin{ex} %[2D1V02]
Tìm số liền sau của số liền sau số $48$.
\shortans{50}
\loigiai{
\textbf{Đáp số:} $50$\\
$48 + 1 + 1 = 50$.
}
\end{ex}

\begin{ex} %[2D1V02]
Tìm số liền trước của số liền trước số $62$.
\shortans{60}
\loigiai{
\textbf{Đáp số:} $60$\\
$62 - 1 - 1 = 60$.
}
\end{ex}

\begin{ex} %[2D1V02]
Có bao nhiêu số tự nhiên nằm giữa số $40$ và số $45$?
\shortans{4}
\loigiai{
\textbf{Đáp số:} $4$\\
Các số là $41, 42, 43, 44$ (có 4 số).
}
\end{ex}

\begin{ex} %[2D1C02]
Tìm một số biết rằng số liền sau của nó là số bé nhất có hai chữ số giống nhau.
\shortans{10}
\loigiai{
\textbf{Đáp số:} $10$\\
Số bé nhất có 2 chữ số giống nhau là 11. Số cần tìm là $11 - 1 = 10$.
}
\end{ex}
